\documentclass[a4paper]{article}

\usepackage[a4paper, top=8mm, bottom=8mm, left=8mm, right=8mm]{geometry}

\usepackage{polyglossia}
\setdefaultlanguage[babelshorthands=true]{russian}
\setotherlanguage{english}

\usepackage{fontspec}
\setmainfont{FreeSerif}
\newfontfamily{\russianfonttt}[Scale=0.7]{DejaVuSansMono}

\usepackage[tiny, compact]{titlesec}

\usepackage{titling}
\setlength{\droptitle}{-1cm}
\setlength{\parskip}{1.5ex}
\pretitle{\begin{center}\begin{bfseries}\Large}
\posttitle{\par\end{bfseries}\end{center}}
\preauthor{\begin{center}\normalsize}
\postauthor{\par\end{center}\vspace{-1.8cm}}

\usepackage{enumitem}
\setlist[enumerate]{topsep=0pt, itemsep=0.2pt, leftmargin=2em}

\usepackage{hyperref}
\usepackage{bookmark}
\usepackage{csquotes}
\usepackage{amsmath}
\usepackage{graphicx}

\title{Реализация свёртки изображений}

\author{Парфёнова Виктория Вадимовна}

\date{}

\begin{document}

\maketitle

\begin{flushright}
    Группа: \emph{25.Б42-мм}

    Кафедра: \emph{системного программирования СПбГУ}

    Научный руководитель: \emph{Пономарев Николай Алексеевич}
    
    Номер семестра практики: \emph{2}
\end{flushright} 

\section{Цель работы}

Целью работы является изучение алгоритма двумерной свёртки, его программная реализация и сравнение производительности полученной реализации с реализацией из библиотеки OpenCV.

\section{Теория}

В ходе работы были изучены следующие материалы по данной теме:

\begin{enumerate}
    \item Image filtering using convolution in OpenCV (документация OpenCV)\footnote{Image filtering using convolution in OpenCV --- URL: \url{https://opencv.org/blog/image-filtering-using-convolution-in-opencv/} (дата обращения: 20.05.2026).}.
    \item Convolutions on Images (Algorithm Archive)\footnote{Convolutions on Images --- URL: \url{https://www.algorithm-archive.org/contents/convolutions/2d/2d.html} (дата обращения: 20.05.2026).}.
\end{enumerate}

Свёртка заключается в следующем: по изображению скользит ядро размером kxk с шагом 1. Для каждой позиции вычисляется сумма поэлементных произведений пикселей окна и коэффициентов ядра:

\[
        O[i,j] = \sum_{u=0}^{k-1} \sum_{v=0}^{k-1} I_{\text{pad}}[i+u, j+v] \cdot K[u,v]
    \]

    \begin{tabular}{l}
        где: \\
        $O[i,j]$ --- значение для пикселя нового изображения; \\
        $I_{\text{pad}}[i+u, j+v]$ --- матрица размером с ядро. Центр матрицы --- исходный пиксель; \\
        $K[u,v]$ --- ядро.
    \end{tabular}

Ядро свёртки (фильтр) --- это небольшая, обычно квадратная матрица коэффициентов, которая определяет эффект преобразования изображения. В процессе свёртки ядро последовательно сдвигается по изображению, и для каждого положения вычисляется взвешенная сумма значений пикселей под ядром.

Пример ядра (sharpen\_3x3 --- повышение резкости):

\[
\begin{pmatrix}
0 & -1 & 0 \\
-1 & 5 & -1 \\
0 & -1 & 0
\end{pmatrix}
\]

Краевые пиксели либо обрезаются (no\_padding), либо обрабатываются дополнением нулями (zero\_padding), дублированием краевых пикселей (replicate\_padding) и др., что позволяет сохранить размер выходного изображения равным входному.

Обработка краёв необходима, потому что у граничных пикселей изображения нет полного набора соседей (например, у левого верхнего угла нет соседей слева и сверху). Если ничего не предпринимать, выходное изображение будет меньше исходного на размер ядра. Чтобы сохранить размерность, применяется дополнение (padding): добавляются искусственные пиксели (нули, отражённые, скопированные и т.п.) перед свёрткой.

\section{Реализация}

Для реализации был выбран язык программирования Python 3.10, так как он предоставляет широкий набор библиотек (NumPy, Pillow, Matplotlib и др.), необходимых для реализации свёртки, визуализации результатов и сравнения производительности. Кроме того, раннее знакомство с данным языком програмирования позволило сосредоточиться на достижении поставленой цели, а не на изучении синтаксиса.

В рамках работы реализованы функции convolution\_grayscale и convolution\_rgb, которые принимают изображение (в виде массива NumPy), ядро свёртки и режим обработки краёв.

Алгоритм работы:

\begin{enumerate}
    \item Исходное изображение дополняется или обрезается согласно выбранному режиму обработки краёв.
    \item Выполняется свёртка (для RGB-изображений операция выполняется независимо для каждого канала, после чего каналы объединяются).
    \item Функция возвращает итоговое изображение (в виде массива NumPy).
\end{enumerate}

Выбор конкретного ядра, режима обработки краёв, входного изображения, выходной директории, количества каналов (grayscale или RGB-изображение) осуществляется через аргументы командной строки с использованием библиотеки Typer\footnote{Typer --- URL: \url{https://typer.tiangolo.com/} (дата обращения: 20.05.2026).}.

Для проверки корректности собственной реализации свёртки использовалась функция scipy.signal.convolve2d, так как именно она позволяет проверять режим no\_padding (обрезка краёв, соответствует valid)\footnote{SciPy API --- Signal processing (scipy.signal) --- convolve2d --- URL:\url{https://docs.scipy.org/doc/scipy/reference/generated/scipy.signal.convolve2d.html} (дата обращения: 21.05.2026).}. 
cv2.filter2D из OpenCV всегда возвращает изображение того же размера, что и входное\footnote{Function opencv::imgproc::filter\_2d --- URL:\url{https://docs.opencv.org/4.x/dd/d6a/tutorial_js_filtering.html}\label{opencv} (дата обращения: 21.05.2026).}, поэтому не позволяет проверить все варианты обработки краёв, реализованные в работе.

Реализация решения осуществлялась с использованием также следующего программного стека:

\begin{enumerate}
    \item NumPy --- для работы с операциями над двумерными массивами.
    \item Pillow --- для загрузки и сохранения изображений, преобразования в массивы NumPy.
    \item Matplotlib --- для визуализации результатов и построения графиков производительности.
    \item pytest --- для тестирования.
    \item pytest-benchmark --- для замера производительности и сравнения времени выполнения.
\end{enumerate}

Управление зависимостями выполнялось через инструмент uv\footnote{uv --- URL: \url{https://docs.astral.sh/uv/} (дата обращения: 20.05.2026).}.

Доступны следующие ядра: 

\begin{enumerate}
    \item sharpen\_3x3: повышение резкости.
    \item blur\_3x3: усредняющее размытие.
    \item gaussian\_blur\_3x3: размытие по Гауссу.
    \item highlighting\_vertical\_borders\_3x3: выделение вертикальных границ.
    \item highlighting\_horizontal\_borders\_3x3: выделение горизонтальных границ.
    \item embossing\_3x3: эффект тиснения.
    \item blur\_5x5: усредняющее размытие (более плавное размытие, чем у 3х3).
    \item gaussian\_blur\_5x5: размытие по Гауссу (более плавное размытие, чем у 3х3).
\end{enumerate}

Доступны следующие варианты обработки края:

\begin{enumerate}
    \item no\_padding: игнорирование краёв (результат будет меньше исходного изображения на размер ядра).
    \item zero\_padding: дополнение нулями.
    \item replicate\_padding: копирование крайних пикселей.
    \item mirror\_padding: отражение без дублирования границы.
    \item symmetric\_padding: симметричное отражение с дублированием границы.
    \item esymmetric\_padding`: симметричное отражение с дублированием границы.
    \item tile\_padding: циклическое продолжение.
\end{enumerate}

Репозиторий с реализацией: \url{https://github.com/roriess/image-convolution}

\section{Эксперименты}

Проведён эксперимент по сравнению скорости выполнения свёртки между учебной реализацией (на языке Python с использованием вложенных циклов и NumPy) и библиотечной функцией OpenCV (cv2.filter2D). Реализация OpenCV написана на C++ и содержит низкоуровневые оптимизации, например, используется алгоритм с SIMD-инструкциями (команды, которые выполняют одну операцию сразу над несколькими элементами данных)\ref{opencv}.

Измерения проводились для квадратных изображений трёх размеров: 128х128, 512х512 и 2048х2048 пикселей. Для каждого размера использовались несколько ядер и режимов обработки краёв.

\begin{figure}[h!]
    \centering
    \includegraphics[width=0.9\textwidth]{images/benchmark\_results.png}
    \caption{Сравнение производительности библиотеки OpenCV и учебной реализации}
    \label{fig:benchmark_results}
\end{figure}

Вывод, полученный из экспериментов: pеализованная свертка уступает в производительности функции cv2.filter2D  из OpenCV\footnote{Benchmark results description --- URL: \url{https://github.com/roriess/image-convolution/blob/main/src/benchmark/benchmark_results_description.md} (дата обращения: 20.05.2026).}.;

\section{Заключение}

В ходе выполнения работы были достигнуты следующие цели:

\begin{enumerate}
    \item Изучена теория по свертке и необходимым для практической части пакетам.
    \item Реализована свертка для grayscale и RGB изображений с возможным выбором ядра и варианта обработки края. 
    \item Поставлены эксперименты.
\end{enumerate}

\end{document}